\section{The \instname\ data reduction flow.}
\label{sec:reduction_flow}


The overall data flow of the \instname\ pipeline is displayed in Figure \ref{fig:reduction_cascade}.

 
\begin{figure}[h!]
\centering
\includegraphics[width=0.9\textwidth]{../instrument_tutorial/figures/reduction_cascade.jpg}
\caption{The data reduction cascade of the \instname\ workflow.}
\label{fig:reduction_cascade}
\end{figure}



The reduction cascade is organized in tasks, which represent
well-defined steps in the process. Tasks can be grouped inside
sub-workflows.  Each task runs a recipe; the detailed description of
the algorithms, input, outputs and recipe parameters used in each
recipe are available in the pipeline manual. Here, we present only the
description of most important features.

The \wkfname\  \edps\ workflow is designed to execute the tasks that
deliver the final reduced data cube for each dataset. It can be either
the product of a single exposure, or the combination of multiple
exposures. Only calibrations needed by the selected the scientific
exposures are processed.

It is possible to set \edps\ to perform the data reduction until a
certain step of the reduction chain (e.g. to reduce only standar
stars, or only flat fields).  This is done by specifying the desired
tasks in the field {\bf Select reduction target} of the {\bf Raw Data}
tab.

The reduction steps of the \wkfname\ workflow are listed below. Before
starting the reduction, the parameters of the recipes associated to
each task can be configured by pressing the button

\includegraphics[width=0.8cm]{../common/figures/configure_dataset.jpg}
close to each dataset configuration.
See for more info on the configuration editor \ref{sec:configuration_editor}



% {\bf LIST THE REDUCTION STEPS HERE}

The \instname\ data reductions is closely intertwined with the observing mode.
Typically, the \instname\ observations are carried out in the so-called nodding mode, in which the telescope is ``nodding'' so the target is placed at two positions along the slit -- nodding positions A and B. In this case, in sequential frames the instrument sees sky at the position where the target was located in the preceding or the following exposure. This allows to remove the sky contribution to zeroth order just by taking the difference between frames taken at two different nodding positions. Brighter targets that require shorter exposures, thereby minimizing the sky emission, may be observed in stare mode -- in this case the telescope remains stationary, so the target is always kept on the same position on the slit. In this case, if there is any significant sky emission, it must be extracted from off-target apertures, and then subtracted from the target aperture. Finally, large targets that occupy more than half of the slit are typically observed in offset mode -- the telescope alternates between target exposures and clear sky exposures. Applying small offsets along the slit, called jittering, may be introduced in all these cases. Jitters help to minimize the effect of poor detector cosmetics, because the same point of the target is observed with different pixels in different exposures. For further details, see \instname\ Pipeline Manual.


\subsection{Detector linearity }

The detector's response is non-linear. In other words the measured ADUs are not linearly proportional to the number of incident photons and the correction is different for each pixel.
The correction is applied on a calibration data product (PRO.CATG=CAL\_DETLIN\_COEFFS) provided by the ESO Archive. The correction is applied by the {\tt cr2res\_cal\_detlin} recipe. At $\sim$30,000\,ADU, the non-linearity is at $\sim$10\,\% level, and the correction is accurate to 0.5\,\%. During the derivation of the non-linearity correction bad pixels that do not follow the expected linearity behavior are identified and masked out. Further details are available in the \instname\ Pipeline manual, Sec.\,6.2.5. 


\subsection{Dark}

Infrared detectors are sensitive to the thermal electrons existing inside the detector itself, and despite the cooling, they always exhibit some intrinsic signal that is proportional to the integration's length. The dark frames can show structures because of the different characteristics of the ADCs and heat sources inside the instrument (e.g., amplifier glow). To remove the dark signal and associated structures, a master dark is subtracted from each frame. It is true that in some cases the dark can be subtracted together with the sky frames -- for example in nodding mode observations. However, subtracting the dark in advance is often preferable, because it improves the statistic calculations necessary to perform rejections by removing the patterns discussed above. The master dark is generated by the {\tt cr2res\_cal\_dark} recipe and is subtracted in various science recipes using the utility recipe {\tt cr2res\_util\_calib}. The dark signal in the flats is removed, typically by subtracting the lamp OFF frames from the lamp ON frames. Note that removing the dark signal does not remove the noise contribution that comes from it. The dark frames are used to identify a type of bad pixels called ``hot'' -- that always exhibit high values. Further information is available in the \instname\ Pipeline manual, Sec.\,4.4. 


\subsection{Flat}

Even the best detectors are not uniform, and individual pixels show slightly different sensitivity to light. The usual approach to correct for this effect is to observe a uniformly illuminated source and derive the true sensitivity of each pixel. Note that the pixels may suffer from intra-pixel sensitivity as well, which cannot be corrected by means of flat fielding.
The master flat is generated by the {\tt cr2res\_cal\_flat} and it is applied by various science recipes using the utility recipe {\tt cr2res\_util\_calib}. The illuminated flat frames are used to identify a type of bad pixels called ``dead'' -- that always exhibit low, near-zero values. This step is also used to identify the location of the spectral orders. The result is stored in a tracewave table with geometric information about the orders. Further information is available in the \instname\ Pipeline manual, Sec.\,4.4.


\subsection{Wavelength calibration}

The \instname\ wavelength calibration depends on the mode, particularly on the observed wavelength range. The user can control it with the {\tt ---wl\_method} parameter of the {\tt cr2res\_util\_wave} recipe. Generally, the calibration is derived by cross-
correlating the extracted lamp spectrum and a synthetic Uranium--Neon spectrum. For the Fabry--P\'erot, the wavelength calibration is derived with a fit to the 2D spectrum, but this method still needs the UNe lamp calibration as an initial guess, because the Fabry--P\'erot spectra show regular, repeatable patterns. The individual calibration strategies are implemented in the {\tt cr2res\_util\_wave} recipe, and their combination is implemented in {\tt cr2res\_cal\_wave} recipe. These recipes update the tracewave table derived in the previous step with the wavelength information. Further information is available in the \instname\ Pipeline manual, Sec.\,6.2.3.
Note that for gas cell wavelength calibration (that only works for the shorter wavelengths) a special external tool\footnote{VIPER (Velocity and IP EstimatoR): \url{https://mzechmeister.github.io/viper_RV_pipeline/}} is available for improving 
the accuracy, and it has been demonstrated to improve the precision down to 2-3 m\,s$^{-1}$.

\subsection{Standard star}

The instrument response is derived by observing a set of flux-calibrated standard stars. Their true spectra are stored in a FITS table {\tt std\_photom.fits} that has self-explanatory column names and is located in the calibration sub-directory of the directory where the pipeline is installed: {\tt .../install/calib/cr2re-1.6.11/}. The pipeline derives the response function by comparing the shape of the observed spectrum with the shape of the supplied flux calibrated spectrum in that table. 

An example of full manual command-line sequence how to calibrate the \instname\ data is available in the \instname\ Pipeline Manual, Sec.\,5.5. Inspecting it can be helpful for users, because EDPS follows the same strategy, even though it is executed automatically.


\subsection{Science}

This step is observing mode and wavelength dependent, and it is executed by one of the following recipes: {\tt cr2res\_obs\_staring}, {\tt cr2res\_obs\_nodding}, {\tt cr2res\_obs\_2d}, and {\tt cr2res\_obs\_pol}. These recipes apply all calibrations listed above and (except {\tt cr2res\_obs\_2d}) perform extraction of science targets. Both regular aperture and optimal (recommended) extraction methods are available. That step is the most likely to need refinement and the pipeline users will benefit from experimenting with the extraction parameters. These include: 
\begin{itemize}
\item {\tt ---extract\_height} that defines the extraction height from the edge-tracing polynomials, defined in the tracewave file; this is particularly important if there are multiple sources on the slit or if the target is not well centered. 
\item {\tt ---extract\_swath\_width} that defines the length of the regions in which, among other things, the target profile for the optimal extraction is defined; extracting spectra of fainter targets benefits from higher values of this parameters.
\item {\tt ---extract\_oversample}, which controls the oversampling of the target profile; experiments suggest that values $>$12 work well.
\end{itemize}
Detailed information about the optimal extraction can be found in the \instname\ Pipeline Manual, Sec.\,6.1.1. Note that the optimal extraction does not work for targets with variable spatial profiles, for example resolved galaxies and disks that may have different profiles for the continuum and for various emission lines.
